现在设 $\Omega$ 为 $\mathbb{R}^m$ 中具有光滑边界的有界区域.\footnote{由于本节着重介绍基本概念, 而不精确陈述结果, 我们对光滑性假设将不作严格规定. 因而, ``光滑'' 应理解为 ``按需要足够光滑''.} 设 $L(\bm{x},D)$ 为定义在 $\overline\Omega$ 上, 具有光滑系数的 $2p$ 阶微分算子. 此外, 设 $B_j(\bm{x},D)$, $j=1,\ldots,p$, 为对 $\bm{x}\in\partial\Omega$ 定义的阶数低于 $2p$ 的微分算子. 以下关注边值问题
\begin{equation}\label{eq:general-homogeneous-boundary-value-problem}
\begin{aligned}
L(\bm{x},D)u&=f(\bm{x}),\quad\bm{x}\in\Omega,\\
B_j(\bm{x},D)u&=0,\quad\bm{x}\in\partial\Omega,\quad j=1,\ldots,p.
\end{aligned}
\end{equation}


假设还存在对 $\bm{x}\in\partial\Omega$ 定义的微分算子 \[S_j(\bm{x},D),\quad T_j(\bm{x},D),\quad C_j(\bm{x},D),\quad j=1,\ldots,p,\] 具有以下性质:\begin{enumerate}
\item $S_j$, $T_j$ 和 $C_j$ 具有光滑系数, 且阶数低于 $2p$.
\item 对任意一组定义在 $\partial\Omega$ 上的光滑函数 $\phi_j$, $j=1,\ldots,2p$, 存在函数 $u,v\in C^{2p}(\overline\Omega)$, 使在 $\partial\Omega$ 上有 $B_ju=\phi_j$, $S_ju=\phi_{p+j}$, $j=1,\ldots,p$, 以及分别有 $C_jv=\phi_j$, $T_jv=\phi_{p+j}$, $j=1,\ldots,p$.
\item 对任意 $u,v\in C^{2p}(\overline\Omega)$, 有
\end{enumerate}
\begin{equation}\label{eq:general-boundary-operator-green-identity}
\begin{aligned}
&\int_\Omega vL(\bm{x},D)u-uL^*(\bm{x},D)v\,\mathrm{d}\bm{x}\\
&\quad=\int_{\partial\Omega}\sum_{j=1}^p S_j(\bm{x},D)uC_j(\bm{x},D)v
-B_j(\bm{x},D)uT_j(\bm{x},D)v\,\mathrm{d}S.
\end{aligned}
\end{equation}


当然, 边值问题必须满足什么假设才能使这样的算子存在, 是极为重要的问题; 稍后讨论椭圆边值问题时将回到这一问题. 目前, 只把 $S_j$, $T_j$ 和 $C_j$ 的存在视为已知.

\begin{Definition}
\label{def:adjoint-boundary-value-problem}
设前述假设成立. 那么边值问题
\begin{equation}\label{eq:general-adjoint-boundary-value-problem}
\begin{aligned}
L^*(\bm{x},D)v&=g(\bm{x}),\quad\bm{x}\in\Omega,\\
C_j(\bm{x},D)v&=0,\quad\bm{x}\in\partial\Omega,\quad j=1,\ldots,p.
\end{aligned}
\end{equation}
称为 \eqref{eq:general-homogeneous-boundary-value-problem} 的\textbf{伴随}边值问题.
\end{Definition}

注意, 若 $u$ 和 $v$ 分别满足 \eqref{eq:general-homogeneous-boundary-value-problem} 和 \eqref{eq:general-adjoint-boundary-value-problem}, 则根据 \eqref{eq:general-boundary-operator-green-identity}, 有
\begin{equation}\label{eq:adjoint-solutions-forcing-orthogonality}
\int_\Omega fv-gu\,\mathrm{d}\bm{x}=0.
\end{equation}
我们并未声称算子 $C_j$ 唯一; 事实上, 即使对于常微分方程, 伴随边界条件也只在线性重组的意义下确定. 不过, 可以对条件 $C_jv=0$ 所刻画的函数集合给出内在刻画.

\begin{Lemma}
\label{lem:adjoint-boundary-intrinsic-characterization}
设 $v\in C^{2p}(\overline\Omega)$, 并用 $X_B$ 表示所有在 $\partial\Omega$ 上对 $j=1,\ldots,p$ 满足 $B_ju=0$ 的 $u\in C^{2p}(\overline\Omega)$ 组成的集合. 那么, 对每个 $u\in X_B$ 都有 $(v,Lu)=(L^*v,u)$ 的充分必要条件是对 $j=1,\ldots,p$ 有 $C_jv=0$.
\end{Lemma}

\begin{Proof}
一个方向由 \eqref{eq:general-boundary-operator-green-identity} 显然成立. 为证明反向结论, 注意由假设可以构造 $u\in C^{2p}(\overline\Omega)$, 使 $B_ju=0$ 且 $S_ju=\phi_j$, 其中 $\phi_j$ 为给定的光滑函数. 若 $(v,Lu)=(L^*v,u)$, 则 \eqref{eq:general-boundary-operator-green-identity} 化为
\begin{equation}\label{eq:adjoint-boundary-conditions-test-identity}
\int_{\partial\Omega}\sum_{j=1}^p\phi_j(\bm{x})C_j(\bm{x},D)v\,\mathrm{d}S=0.
\end{equation}
若这对任意 $\phi_j$ 成立, 则 $C_jv$ 显然必为零.
\end{Proof}

因此, 尽管可能存在不同的伴随边界条件组, 它们必定彼此等价. 需要提醒的是, 等价的边界条件组不必互为线性组合. 例如, 设 $\partial\Omega$ 为 $\mathbb{R}^2$ 中的闭曲线, 并用 $s$ 表示弧长. 那么条件 $v=0$ 和 $\mathrm{d}v/\mathrm{d}s+v=0$ 等价, 尽管它们并不互为倍数.

以两个例子结束本小节:

\Needspace{4\baselineskip}
\begin{Example}
\label{ex:laplace-dirichlet-neumann-self-adjoint}
有
\begin{equation}\label{eq:laplace-dirichlet-neumann-green-identity}
\int_\Omega u\Delta v-v\Delta u\,\mathrm{d}\bm{x}=\int_{\partial\Omega}u\dfrac{\partial v}{\partial n}-v\dfrac{\partial u}{\partial n}\,\mathrm{d}S.
\end{equation}
因此 Laplace 方程的 Dirichlet 和 Neumann 边值问题都是它们自身的伴随.
\end{Example}

\Needspace{4\baselineskip}
\begin{Example}
\label{ex:normal-tangential-adjoint-boundary-problem}
设 $\Omega$ 为 $\mathbb{R}^2$ 中以闭光滑曲线为边界的有界区域. 用 $s$ 表示沿该曲线的弧长. 考虑边值问题
\begin{equation}\label{eq:laplace-normal-tangential-boundary-problem}
\Delta u=f(\bm{x}),\ \bm{x}\in\Omega,\qquad\dfrac{\partial u}{\partial n}+\dfrac{\partial u}{\partial s}=0,\ \bm{x}\in\partial\Omega.
\end{equation}
得到
\begin{equation}\label{eq:laplace-normal-tangential-adjoint-calculation}
\begin{aligned}
\int_\Omega u\Delta v-v\Delta u\,\mathrm{d}\bm{x}
&=\int_{\partial\Omega}u\dfrac{\partial v}{\partial n}-v\dfrac{\partial u}{\partial n}\,\mathrm{d}s\\
&=\int_{\partial\Omega}u\left(\dfrac{\partial v}{\partial n}-\dfrac{\partial v}{\partial s}\right)
-v\left(\dfrac{\partial u}{\partial n}+\dfrac{\partial u}{\partial s}\right)\,\mathrm{d}s.
\end{aligned}
\end{equation}
因此伴随边值问题为
\begin{equation}\label{eq:laplace-normal-tangential-adjoint-problem}
\Delta v=g(\bm{x}),\ \bm{x}\in\Omega,\qquad\dfrac{\partial v}{\partial n}-\dfrac{\partial v}{\partial s}=0,\ \bm{x}\in\partial\Omega.
\end{equation}

\end{Example}
